\chapter{Implementação}
\label{cap:implem}

% figuras estão no subdiretório "figuras/" dentro deste capítulo
\graphicspath{\currfiledir/figuras/}

A implementação do Path Tracer foi feita a partir de um Ray Tracer básico,
implementado seguindo o texto de \citet{shirley25}. A partir dele foram feitas
modificações principalmente para incorporar as diferentes técnicas de
amostragem e a combinação delas por meio do amostragem de importância múltipla.
Além disso foram implementadas algumas imagens para visualizar a contribuição
de cada método e a vâriancia em cada pixel, além de outros materiais e suas
correspondentes distribuições de probabilidade.

A saída do programa são três imagens:

\begin{description}
  \item[output.png] A imagem gerada pelo algoritmo de Path Tracing.
  \item[variance.png] Uma imagem gerada com valores de variância de cada pixel,
    permitindo achar locais onde podemos melhorar nosso algoritmo.
  \item[mis\_viz] Uma imagem de visualização das técnicas de amostragem, onde
    cada pixel é colorido conforme a técniaca usada.
\end{description}

O código foi produzido na linguagem \emph{Rust}, usando o compilador \texttt{rustc},
versão \texttt{1.96.0}.

\section{Arquitetura}

Seguindo \citet{shirley25} temos alguns módulos, cada qual em um arquivo. Eles
são

\begin{description}
  \item[main.rs] Ponto de entrada do programa, é onde está o
código de \emph{setup} inicial, assim como a descrição da cena e a função
principal que irá estimar a equação da renderização.
  \item[utils.rs] Módulo de utilidades, possui funções para obter números
    aleatórios e algumas funções de operações comuns em vetores.
  \item[vec3.rs] Possui a definição da estrutura de dados \texttt{Vec3} e
    seus métodos, que serão usados amplamente no programa.
  \item[hittable.rs] Define a interface \texttt{Hittable}, que abstrai um
    elemento da cena que pode ser atingido por um raio de luz. Além disso
    também define o \texttt{HittableList}, uma coleção de elementos
    \texttt{Hittable}.
  \item[sphere.rs] Define a estrutura de dados \texttt{Sphere}, que representa
    uma esfera na cena. Além disso implementa a interface \texttt{Hittable}
    para ela, permitindo que ela apareça na cena.
  \item[ray.rs] Define a estrutura de dados \texttt{Ray}, que representa um
    raio com um ponto de início e uma direção.
  \item[material.rs] Possui a implementação dos BRDFs utilizados, uma função
    para amostragem baseada em cada BRDF e também a distribuição de
    probabilidade correspondente a cada essa amostragem. Também usamos essa
    interface para fontes de luz, e definimos uma fonte de luz difusa.
  \item[light\_sampler.rs] Possui a definição da estrutura de dados do
    \texttt{LightSampler}, que permitirá amostrar pontos em uma fonte de luz. 
  \item[camera.rs] Define a estrutura \texttt{Camera}, e implementa a função
    que gerará os raios iniciais para cada pixel, baseados nos parâmetros da
    estrutura.
\end{description}


\section{Laço principal}

No arquivo \texttt{main.rs} temos o laço principal do programa, descrito no
Algoritmo \autoref{alg:mainloop}. Nele, iteramos por todos os pixels, para cada um
obtemos um raio inicial com base na câmera, computamos o integrando $N$
vezes usando a função \texttt{L()}, e, finalmente, normalizamos e computamos
valores como a média, a variância e a cor da imagem de visualização.

\begin{algorithm}
\caption{Laço principal}
\label{alg:mainloop}
\begin{algorithmic}[1]

\For{$j = 0$ até altura $-1$}
\For{$i = 0$ até largura $-1$}

    \State \text{média} $\gets 0$
    \State \text{M2} $\gets 0$
    \State \text{mis\_viz} $\gets 0$

    \For{$k = 1$ até $N$}
        \State \text{raio} $\gets \Call{NovoRaio}{camera, i, j}$

        \State $(\text{radiância}, \text{mis}) \gets \Call{L}{\text{raio}, \text{world}, \text{lights}}$

        \State $\text{delta} \gets \text{radiância} - \text{média}$
        \State \text{média} $\gets \text{média} + \text{delta} / k$
        \State \text{M2} $\gets \text{M2} + \text{delta} \cdot (\text{radiância} - \text{média})$
        \State \text{mis\_viz} $\gets \text{mis\_viz} + \text{mis}$

    \EndFor

    \State \text{variância} $\gets \text{M2} / N$
    \State \text{mis\_viz} $\gets \text{mis\_viz} / N$

    \State \text{EscreverPixel}(i, j, \text{média}, \text{variância}, \text{mis\_viz})

\EndFor
\EndFor

\end{algorithmic}
\end{algorithm}

Aqui $média$ é a estimativa de Monte Carlo para a radiância; $variância$ é a
variância empírica de cada pixel, obtida pelo algoritmo de
Welford\footnote{Usamos esse algoritmo para evitar problemas de instabilidade
numérica, já que, calculando usando a fórmula clássica, $\left\langle (X -
\langle X \rangle)^2 \right\rangle$, acabamos fazendo a subtração de dois
números potencialmente muito próximos, o que aumenta significantemente a
magnitude de erros numéricos.} \citep{welford62}; e $mis\_viz$ se refere à um
valor que codifica de qual estratégia de amostragem obtivemos a contribuição
naquele pixel. 

É importante ressaltar que, na implementação, as grandezas radiométricas são
representadas por vetores RGB (\texttt{Vec3}). Dessa forma, os cálculos
descritos anteriormente são executados independentemente para os canais
vermelho, verde e azul. Assim, variáveis como radiância, média, e variância
armazenam três valores escalares, correspondentes às contribuições de cada
canal para a cor final do pixel.

\section{A função $L$}

A função $L$, para estimar valor da radiância na direção dada, fará a soma
ponderada das contribuições dos diversos caminhos amostrados, utilizando a
amostragem de importância múltipla da seção \ref{sec:mis}.

\subsection{Formulação do estimador de importância múltipla}

A ideia é criarmos um caminho base, amostrando cada ponto usando a amostragem
por BRDF. Porém, a cada ponto amostrado, também amostramos as fontes de luz,
criando novos caminhos completos a partir da mesma base. Assim para o caminho
base $(x_1, \ldots, x_{k - 1})$ iremos efetivamente amostrar os caminhos

\begin{gather*}
  (x_1, x_{L,2}), \\
  (x_1, x_2, x_{L,3}), \\
  (x_1, x_2, x_3, x_{L,4}), \\
  \vdots \\
  (x_1, \ldots, x_{k-1}, x_{L,k})
\end{gather*}

onde $x_i$ são os pontos do caminho base e $x_{L,i}$ são os pontos amostrados
de fontes de luz em cada passo. Além desses caminhos, temos que contar também
os casos que, ao amostar por BRDF, amostramos um ponto em uma fonte de luz.
Nesse caso, o caminho base termina precocemente.

Assim, temos duas possíveis maneiras de amostrar cada caminho - uma é com o
último vértice do caminho sendo obtido pela amostragem de fontes de luz; a
outra é com ele sendo amostrado pelo BRDF. Por esse motivo, teremos que
utilizar a amostragem de importância múltipla para combinar as técnicas.

Observe, também, que podemos escrever cada técnica de amostragem de
\emph{caminhos} em uma tupla $(k, T)$, com $k \in \mathbb{N}$ e $T \in  \{L,
B\}$, $k$ representando o tamanho do caminho completo e $T$ representando a
técnica de amostragem usada para obter o último vértice do caminho amostrado.

Relembrando a equação \ref{eq:misestimator}, temos

\[
  I = \sum_{i=1}^n w_i(X_{i}) \frac{f(X_i)}{p_i(X_i)}
\]

Para um caminho completo $\bar{x} = (x_1, \ldots, x_k)$, sendo $x_k$ um ponto
em uma fonte de luz, temos

\[
  f(\bar{x}) = L_e(x_{k}) \left( \prod_{i=2}^{k-1} f_s(x_{i-1} \rightarrow
  x_{i} \rightarrow x_{i+1}) \cos \theta_{i+1} \right) 
\]

Para as probabilidades dos caminhos, vamos chamar $p_{k,L}$ a probabilidade de
um caminho ter sido amostrado com seu último vértice por amostragem de fontes
de luz e $p_{k,B}$ a probabilidade com o último vértice utilizando a amostragem
por BRDF. Assim

\begin{align*}
  p_{k,B}(\bar{x}) &= p_B(x_0)p_B(x_1 \mid x_0) \ldots p_B(x_k \mid x_{k-1})\\
  p_{k,L}(\bar{x}) &= p_B(x_0)p_B(x_1 \mid x_0) \ldots p_L(x_k \mid x_{k-1})
\end{align*}

Note que $p_{k,L}$ e $p_{k,B}$ diferem apenas no último termo. Sendo $T \in \{B, L\}$ temos

\[
  w_{k,T}(\bar{x}) = \frac{p_{k,T}(\bar{x})^2}{p_{k,B}(\bar{x})^2 +
  p_{k,L}(\bar{x})^2} = \frac{p_T(x_k \mid x_{k-1})^2}{p_B(x_k \mid x_{k-1})^2
  + p_L(x_k \mid x_{k-1})^2}
\]

e então podemos escrever

\[
  I = \sum_{i=1}^n w_{k,T}(X_{i}) \frac{f(X_i)}{p_{k,T}(X_i)}
\]

\subsection{Implementação}

A função $L$ é implementada como um \emph{loop} infinito, saindo dele em três
ocasiões:

\begin{itemize}
  \item O raio não intersecta nenhuma geometria (ele ``foge da cena'').
  \item O raio atinge um material emissivo. Nesse caso adicionamos a
    contribuição, ponderada corretamente com $w_{\text{k},B}$.
  \item O algoritmo de Roleta Russa decide encerrar o caminho atual.
\end{itemize}

Caso nenhum desses casos não aconteçam, amostramos um ponto nas fontes de luz e
adicionamos sua contribuição com $w_{\text{k}, L}$. Em seguida, amostramos o
BRDF do material atual para obter o próximo ponto do caminho base, atualizamos
as variáveis, e reiniciamos o laço.

A variável \texttt{peso} guardará, na $n$-ésima iteração, o valor correspondente a 

\[
  \frac{ \prod_{i=2}^{n-1} f_s(x_{i-1} \rightarrow x_{i} \rightarrow x_{i+1})
  \cos \theta_{i+1}}{p_B(x_0)p_B(x_1 \mid x_0) \ldots p_B(x_n \mid x_{n-1})}
\]

que é o valor comum, entre todos os próximos vértices, do termo
$\frac{f(X_i)}{p_{k,T}(X_i)}$ de $I$. Dessa forma, não precisamos guardar os
valores de cada vértice e temos um custo de memória constante.

Os trechos \texttt{AvaliaEmissao}, \texttt{AmostrarLuz} e \texttt{RoletaRussa}
serão desenvolvidos nas próximas seções.

\begin{algorithm}
\caption{Função L}
\begin{algorithmic}[1]
  \Function{L}{raio, mundo, luzes}
  \State $\text{profundidade} \gets 0$
  \State $\text{radiancia} \gets 0$
  \State $\text{peso} \gets 1$
  \Loop
    \State $hit \gets \Call{Intersecta}{\text{raio}, \text{mundo}}$
    \If{$hit$}
      \State $\Call{AvaliaEmissao}{hit.material, \text{profundidade}, \text{peso}, \text{luzes}}$

      \State \Call{AmostrarLuz}{$hit$, peso, mundo, luzes}

      \State \Call{RoletaRussa}{profundidade, peso}

      \State $(\text{proximoRaio}, f_s, p_b) \gets \Call{amostrar}{hit.material}$

      \State $\text{peso} \gets \text{peso} \cdot f_s \cdot \cos \theta / p_b$
      \State $\text{raio} \gets \text{proximoRaio}$
      \State $\text{profundidade} \gets \text{profundidade} + 1$
    \Else
    \State \textbf{break}
    \EndIf
  \EndLoop
  \EndFunction
\end{algorithmic}
\end{algorithm}

\section{\texttt{AvaliaEmissao}}

A função \texttt{AvaliaEmissao} verifica se o ponto atual, amostrado pelo BRDF na
iteração passada, emite radiância. Se não emite, a execução continua sem
alteração. Caso contrário, precisamos adicionar a contribuição à estimativa.

No caso de emitir, existem três possibilidades:

\begin{itemize}
  \item É a primeira iteração. Nesse caso, só há uma maneira de amostrar esse
    caminho, que sai diretamente do sensor até a fonte de luz. Portanto, não é
    utilizado os pesos da amostragem de importância múltipla.
  \item O último material possuia um BRDF que continha uma função $\delta$ de
    Dirac.\footnote{Um exemplo de material que utiliza uma função $\delta$ é o
    BRDF de um espelho perfeito. O valor dele é $1$ apenas na direção de
    perfeita reflexão.} Nesse caso, $p_{k,L}$ será sempre 0 (a probabilidade de
    um ponto amostrado numa fonte de luz ter a direção correta é 0). Assim, a
    única forma de amostrar esse caminho é por meio da amostragem por BRDF e,
    portanto, também não usamos os pesos de amostragem múltipla nesse caso.
  \item No terceiro caso, temos uma amostra que também pode ser gerada pela
    amostragem de fontes de luz. Logo, precisamos utilizar os pesos da
    amostragem múltipla. Para isso, no material emissivo, temos uma função que
    retorna a $p_{k,L}$ daquela fonte de luz. Combinando com $p_{k,B}$ obtido
    da amostra de BRDF da iteração anterior, conseguimos calucular $w_{k,B}$.
\end{itemize}

\begin{algorithm}
\caption{Função AvaliaEmissao}
\begin{algorithmic}[1]
  \Function{AvaliaEmissao}{material, profundidade, peso, luzes}
  \If{material.emissivo}
    \If{profundidade = 0}
      \State $\text{radiancia} \gets \text{radiancia} + \text{material.emitido}$
    \ElsIf{material anterior era delta}
      \State $\text{radiancia} \gets \text{radiancia} + \text{peso} \cdot \text{material.emitido}$
    \Else
      \State $p_L \gets \text{material.pl}$
      \State $w \gets p_B^2 / (p_B^2 + p_L^2)$
      \State $\text{radiancia} \gets \text{radiancia} + w \cdot \text{peso} \cdot \text{material.emitido}$
    \EndIf
    \State \textbf{break}
  \EndIf
  \EndFunction
\end{algorithmic}
\end{algorithm}


\section{\texttt{AmostrarLuz}}

Nessa parte, já sabemos que o material atual não é emissivo, então podemos
amostrar uma fonte de luz para obter um possível caminho completo. Para isso
usamos a estrutura \texttt{LightSampler}, que retorna um ponto na fonte de luz
e a probabilidae $p_L$, na medida de área, de amostra-lo.

Em seguida, verificamos se há alguma obstrução entre o ponto na fonte de luz
amostrado e o ponto atual (que foi amostrado pelo BRDF na iteração anterior).
Se existir uma obstrução, continuamos para o resto do \emph{loop}, sem
adicionar nenhuma contribuição. Caso contrário, verificamos o BRDF para obter
seu valor e o $p_B$ do ponto atual. Além disso, convertemos o $p_L$, que está
em medida de área, para a medida de ângulo sólido. Com isso conseguimos
calcular o peso $w_{k, L}$ e portanto a contribuição para esse caminho.

\begin{algorithm}
\caption{Função AmostrarLuz}
\begin{algorithmic}[1]
  \Function{AmostrarLuz}{hit, peso, mundo, luzes}
  \State $(\text{posicaoLuz}, \text{materialLuz},p_L^A) \gets \Call{amostrar}{\text{luzes}}$
  \State $\text{obstruido} \gets \Call{obstruido}{\text{hit.pos}, \text{posicaoLuz}}$
  \If{\textbf{not} obstruido}
    \State $f_s \gets \Call{BRDF}{\text{hit.material}, \text{posicaoLuz}}$
    \State $p_B \gets \Call{PDF}{\text{hit.material}, \text{posicaoLuz}}$
    \State $p_L \gets \Call{ConverterParaAnguloSolido}{p_L^A, \text{hit}, \text{posicaoLuz}}$
    \State $w \gets p_L^2 / (p_B^2 + p_L^2)$
    \State $\text{radiancia} \gets \text{radiancia} + w \cdot \text{peso} \cdot
      f_s \cdot \cos \theta \cdot \text{materialLuz.emitido} / p_L$
  \EndIf
  \EndFunction
\end{algorithmic}
\end{algorithm}

\section{\texttt{RoletaRussa}}

Para evitar caminhos muitos grandes, visto que quanto maior o caminho menor a
contribuição na imagem final, utilizamos o método de Roleta Russa
\citep{arvo90} para terminar o caminho sem enviesar o estimador.

\begin{algorithm}
\caption{Função RoletaRussa}
\begin{algorithmic}[1]
  \Function{RoletaRussa}{profundidade, peso}
  \If{$\text{profundidade} \ge \text{profundidadeMinima}$}
    \State q $\gets \Call{ComponenteMax}{peso}$
    \If{$\Call{Random}{} \ge \text{q}$}
      \State peso $\gets \text{peso} / \text{q}$
    \EndIf
  \EndIf
  \EndFunction
\end{algorithmic}
\end{algorithm}

\section{Conclusão}

Neste capítulo apresentamos a implementação do integrador de Path Tracing
desenvolvido nesse trabalho. A partir de um Ray Tracer básico foram
incorporadas técnicas de redução de variância, como a amostragem de fontes de
luz (NEE), amostragem por BRDF, sua combinação utilizando amostragem de
importância múltipla e a técnica de Roleta Russa. 

Também foi apresentada a formulação dos caminhos amostrados pelo integrador,
juntamente com suas distribuições de probabilidade, os pesos utilizados pela
amostragem de importância múltipla e as  contribuições de radiância associadas
a cada caminho.

